% GENERATED FILE. Edit canonical lessons, then run npm run generate:books.
% canonical-source-hash: fnv1a64:58a48e55e4276b99
% canonical-lessons: TA-C171-atikkati, TA-C171-cila-nerankalil, TA-C171-eppotum, TA-C171-orupotum, TA-C171-tamatam

\chapter{Often, Sometimes, Always, Never, Delay}
\label{ch:ta-often-sometimes-always-never-delay}
\hlchaptermodality{171}

\begin{chapteropening}
\textbf{By the end of this chapter:} \emph{I can say often, sometimes, always, never and a delay.}

Complete the last lesson of chapter 171: I can say often, sometimes, always, never and a delay.
\end{chapteropening}

\section[\texorpdfstring{\ta{அடிக்கடி} (aṭikkaṭi)}{அடிக்கடி (aṭikkaṭi)}]{\ta{அடிக்கடி} (aṭikkaṭi) — often}

\label{lesson:TA-C171-atikkati}



\begin{quote}
Before the new one: say the Tamil for diarrhoea, then the Tamil for the brain.
\end{quote}

\subsection*{You'll want to know: \ta{அடிக்கடி}}
\textbf{\ta{அடிக்கடி}} — \emph{aṭikkaṭi} — \textquotedblleft{}often\textquotedblright{}.

\begin{grammarlens}[title={What you've built}]
One more word for this chapter.
\end{grammarlens}

\subsection*{Your turn}
\begin{itemize}
\raggedright
  \item \emph{Say it:} \emph{aṭikkaṭi}
  \item \emph{Say it:} \emph{aṭikkaṭi}, once more
  \item \emph{Say it:} say \emph{aṭikkaṭi}
  \item \emph{Recall:} say the Tamil for diarrhoea, then the Tamil for the brain, then say \emph{aṭikkaṭi} again
\end{itemize}

\begin{tcolorbox}[breakable,colback=teal!4,colframe=teal!35!black,title={Before you move on}]
What does \ta{அடிக்கடி} mean? (\textquotedblleft{}Often\textquotedblright{}.) Say it once more.
\end{tcolorbox}
\section[\texorpdfstring{\ta{சில} \ta{நேரங்களில்} (cila nēraṅkaḷil)}{சில நேரங்களில் (cila nēraṅkaḷil)}]{\ta{சில} \ta{நேரங்களில்} (cila nēraṅkaḷil) — sometimes}

\label{lesson:TA-C171-cila-nerankalil}



\begin{quote}
Before the new one: say the Tamil for the brain, then the Tamil for often.
\end{quote}

\subsection*{You'll want to know: \ta{சில} \ta{நேரங்களில்}}
\textbf{\ta{சில} \ta{நேரங்களில்}} — \emph{cila nēraṅkaḷil} — \textquotedblleft{}sometimes\textquotedblright{}.

\begin{grammarlens}[title={What you've built}]
One more word for this chapter.
\end{grammarlens}

\subsection*{Your turn}
\begin{itemize}
\raggedright
  \item \emph{Say it:} \emph{cila nēraṅkaḷil}
  \item \emph{Say it:} \emph{cila nēraṅkaḷil}, once more
  \item \emph{Say it:} say \emph{cila nēraṅkaḷil}
  \item \emph{Recall:} say the Tamil for the brain, then the Tamil for often, then say \emph{cila nēraṅkaḷil} again
\end{itemize}

\begin{tcolorbox}[breakable,colback=teal!4,colframe=teal!35!black,title={Before you move on}]
What does \ta{சில} \ta{நேரங்களில்} mean? (\textquotedblleft{}Sometimes\textquotedblright{}.) Say it once more.
\end{tcolorbox}
\section[\texorpdfstring{\ta{எப்போதும்} (eppōtum)}{எப்போதும் (eppōtum)}]{\ta{எப்போதும்} (eppōtum) — always}

\label{lesson:TA-C171-eppotum}



\begin{quote}
Before the new one: say the Tamil for often, then the Tamil for sometimes.
\end{quote}

\subsection*{You'll want to know: \ta{எப்போதும்}}
\textbf{\ta{எப்போதும்}} — \emph{eppōtum} — \textquotedblleft{}always\textquotedblright{}.

\begin{grammarlens}[title={What you've built}]
One more word for this chapter.
\end{grammarlens}

\subsection*{Your turn}
\begin{itemize}
\raggedright
  \item \emph{Say it:} \emph{eppōtum}
  \item \emph{Say it:} \emph{eppōtum}, once more
  \item \emph{Say it:} say \emph{eppōtum}
  \item \emph{Recall:} say the Tamil for often, then the Tamil for sometimes, then say \emph{eppōtum} again
\end{itemize}

\begin{tcolorbox}[breakable,colback=teal!4,colframe=teal!35!black,title={Before you move on}]
What does \ta{எப்போதும்} mean? (\textquotedblleft{}Always\textquotedblright{}.) Say it once more.
\end{tcolorbox}
\section[\texorpdfstring{\ta{ஒருபோதும்} (orupōtum)}{ஒருபோதும் (orupōtum)}]{\ta{ஒருபோதும்} (orupōtum) — never (with a negative verb)}

\label{lesson:TA-C171-orupotum}



\begin{quote}
Before the new one: say the Tamil for sometimes, then the Tamil for always.
\end{quote}

\subsection*{You'll want to know: \ta{ஒருபோதும்}}
\textbf{\ta{ஒருபோதும்}} — \emph{orupōtum} — \textquotedblleft{}never (with a negative verb)\textquotedblright{}.

\begin{grammarlens}[title={What you've built}]
One more word for this chapter.
\end{grammarlens}

\subsection*{Your turn}
\begin{itemize}
\raggedright
  \item \emph{Say it:} \emph{orupōtum}
  \item \emph{Say it:} \emph{orupōtum}, once more
  \item \emph{Say it:} say \emph{orupōtum}
  \item \emph{Recall:} say the Tamil for sometimes, then the Tamil for always, then say \emph{orupōtum} again
\end{itemize}

\begin{tcolorbox}[breakable,colback=teal!4,colframe=teal!35!black,title={Before you move on}]
What does \ta{ஒருபோதும்} mean? (\textquotedblleft{}Never (with a negative verb)\textquotedblright{}.) Say it once more.
\end{tcolorbox}
\section[\texorpdfstring{\ta{தாமதம்} (tāmatam)}{தாமதம் (tāmatam)}]{\ta{தாமதம்} (tāmatam) — a delay}

\label{lesson:TA-C171-tamatam}



\begin{quote}
Before the new one: say the Tamil for always, then the Tamil for never.
\end{quote}

\subsection*{You'll want to know: \ta{தாமதம்}}
\textbf{\ta{தாமதம்}} — \emph{tāmatam} — \textquotedblleft{}a delay\textquotedblright{}.

\begin{grammarlens}[title={What you've built}]
One more word for this chapter.
\end{grammarlens}

\subsection*{Your turn}
\begin{itemize}
\raggedright
  \item \emph{Say it:} \emph{tāmatam}
  \item \emph{Say it:} \emph{tāmatam}, once more
  \item \emph{Say it:} say \emph{tāmatam}
  \item \emph{Recall:} say the Tamil for always, then the Tamil for never, then say \emph{tāmatam} again
\end{itemize}

\begin{tcolorbox}[breakable,colback=teal!4,colframe=teal!35!black,title={Before you move on}]
What does \ta{தாமதம்} mean? (\textquotedblleft{}A delay\textquotedblright{}.) Say it once more.
\end{tcolorbox}
